\section{St\o rmer's theorem: a complete primitive recursive proof}
\label{sec:stormer}

This section proves the full finiteness-and-enumeration conclusion for consecutive
smooth integers.  The original St\o rmer method, described in the introduction to
Lehmer's paper~\cite{Lehmer1964}, uses coefficients whose prime exponents belong to
$\{0,1,2\}$.  Keeping the exponent $2$, rather than insisting on squarefree
coefficients, is decisive: it makes every prime in the Pell solution's second
coordinate divide the coefficient.  We give the elementary argument and a very
coarse explicit bound on the requisite fundamental solutions.  The bound is chosen
for a transparent proof in finite arithmetic, not for efficient numerical work.

\subsection{Statement and conventions}

Let $S=\{p_1,\ldots,p_k\}$ be a finite set of distinct primes.  A positive integer
is \emph{$S$-smooth} if all its prime factors belong to $S$; thus $1$ is smooth.
By convention $0$ is not smooth.
Put $R=\prod_{p\in S}p$, with the empty product equal to $1$, and define
\begin{equation}\label{eq:stormer-B}
 B(D)=(2D+2)^{\,2(2D+1)^3+2}\qquad(D\ge1).
\end{equation}
All exponents in this formula are natural-number exponents.

\begin{theorem}[Primitive recursive St\o rmer theorem]
\label{thm:stormer-pr}
There is a primitive recursive procedure which, from a finite list of distinct
primes $S$, returns precisely all positive integers $n$ for which $n$ and $n+1$
are both $S$-smooth.  Its output has at most $3^k-2^k$ entries.  Every output
satisfies
\[
             2n+1<B(R^2).
\]
If $2\notin S$, the output is empty.
\end{theorem}

We will explicitly prove every ingredient needed for this statement.  Square-root
notation in the proof is shorthand for exact arithmetic with pairs
$a+b\sqrt D$.  For nonsquare natural $D$, the sign of such a pair with rational
coefficients is decided by the signs of $a,b$ and a comparison of $a^2$ with
$Db^2$.  Alternatively, all occurrences of $\lfloor\sqrt m\rfloor$ are the
bounded-search function
\[
 \operatorname{isqrt}(m)=\max\{r\le m:r^2\le m\}.
\]
There is no real-number oracle or appeal to completeness.

\subsection{A bounded Pell construction}

\begin{lemma}[Finite approximation of a quadratic radical]
\label{lem:stormer-dirichlet}
For nonsquare $D\ge2$ and natural $N\ge1$, bounded search constructs positive
integers $p,q$, with $1\le q\le N$, such that
\[
                 |p-q\sqrt D|<1/N.
\]
\end{lemma}

\begin{proof}
For $0\le h\le N$ put
\[
 a_h=\operatorname{isqrt}(Dh^2),\qquad
 b_h=\operatorname{isqrt}(DN^2h^2)-Na_h.
\]
These integers satisfy
\[
 a_h\le h\sqrt D<a_h+1,
 \qquad
 b_h\le N(h\sqrt D-a_h)<b_h+1,
\]
so $0\le b_h<N$.  Among $N+1$ entries taking $N$ possible values, bounded
search finds $i<j\le N$ with $b_i=b_j$.  Set $q=j-i$ and $p=a_j-a_i$.
The difference of the two fractional parts has absolute value less than $1/N$,
which gives the required approximation.  Since $\sqrt D>1$ and $q\ge1$, the
integer $p$ is positive.  Finite pigeonhole here is the elementary counting fact
for bounded lists, proved by induction on the number of boxes.
\end{proof}

\begin{lemma}[An explicit bounded Pell solution]
\label{lem:stormer-pell-bound}
If $D\ge2$ is nonsquare, then bounded search finds positive integers $x,y$ such
that
\[
 x^2-Dy^2=1,
 \qquad x<B(D),\quad y<B(D).
\]
\end{lemma}

\begin{proof}
Write
\[
 C=2D+1,\quad L=C+1,\quad K=2C^3,\quad N_j=L^{j+1}
                    \quad(0\le j\le K).
\]
Apply Lemma~\ref{lem:stormer-dirichlet} with $N=N_j$ and denote its output by
$(p_j,q_j)$.  Put
\[
 t_j=p_j^2-Dq_j^2,\qquad U_j=p_j+q_j\sqrt D.
\]
The integer $t_j$ is nonzero: otherwise $D$ would be a square, by unique
factorization applied to $p_j^2=Dq_j^2$.  Moreover
\[
 U_j<2Dq_j+1\le Cq_j,
 \qquad
 0<|t_j|=|p_j-q_j\sqrt D|\,U_j<C.
\]
The displayed strict upper estimate for $U_j$ follows from $\sqrt D<D$ and
$|p_j-q_j\sqrt D|<1/N_j\le1$.
Because $|t_j|\ge1$, the same estimates give
\[
 \frac1{Cq_j}<\frac1{N_j},\qquad q_j>N_j/C.
\]
Consequently $q_{j+1}>N_{j+1}/C>N_j\ge q_j$.

Color the $K+1$ pairs by
\[
       \bigl(t_j,\ p_j\bmod |t_j|,\ q_j\bmod |t_j|\bigr).
\]
There are at most
$2\sum_{t=1}^{C-1}t^2\le2C^3=K$ colors.  A second finite pigeonhole search
therefore gives $i<j$ with the same color.  Write their common nonzero norm as
$t$.  Define
\[
 x=\frac{p_jp_i-Dq_jq_i}{t},\qquad
 y=\frac{q_jp_i-p_jq_i}{t}.
\]
Both quotients are integers: the first numerator is congruent modulo $|t|$ to
$p_i^2-Dq_i^2=t$, and the second is congruent to zero.  Multiplication in the
quadratic ring gives
\[
 x+y\sqrt D=\frac{U_j}{U_i},\qquad x^2-Dy^2=1.
\]
Since $q_j>q_i$ and the norms agree, $p_j>p_i$, so $U_j>U_i>1$.
For a norm-one number $V=U_j/U_i>1$,
\[
 x=\frac{V+V^{-1}}2>1,\qquad
 y=\frac{V-V^{-1}}{2\sqrt D}>0.
\]
Thus $x,y$ are positive integers, and
\[
 x<V<U_j<CN_j\le C L^{K+1}<L^{K+2}=B(D),
 \qquad y<V<B(D).
\]
All searches just specified have explicit natural-number bounds.
\end{proof}

\subsection{The structure of the positive Pell solutions}

Let $(u,v)$ be the positive solution of $u^2-Dv^2=1$ with least $u$; it can be
found by searching $1\le y<B(D)$, setting $x=\operatorname{isqrt}(Dy^2+1)$,
and taking the first exact solution.  Minimizing $y$ is equivalent to minimizing
$x$.  Set $\varepsilon=u+v\sqrt D$ and define integer pairs recursively by
\begin{align*}
 x_0&=1,&y_0&=0,\\
 x_{m+1}&=ux_m+Dvy_m,&y_{m+1}&=vx_m+uy_m.
\end{align*}
Then $x_m+y_m\sqrt D=\varepsilon^m$, and each pair has norm one.

\begin{lemma}[Bounded generation by the fundamental solution]
\label{lem:stormer-pell-powers}
Every positive solution $(x,y)$ of $x^2-Dy^2=1$ equals $(x_m,y_m)$ for a unique
positive index $m\le x$.
\end{lemma}

\begin{proof}
The coordinates of positive norm-one solutions are ordered together:
$x'<x$ if and only if $y'<y$, and then
$x'+y'\sqrt D<x+y\sqrt D$.  Since $u\ge2$, the recursion shows that
$x_m\ge m+1$ and that $x_m$ is strictly increasing.  Bounded search therefore
finds the largest $m\le x$ such that $x_m\le x$; necessarily $m<x$, and
\[
     \varepsilon^m\le x+y\sqrt D<\varepsilon^{m+1}.
\]
The quotient
\[
 \gamma=(x+y\sqrt D)(x_m-y_m\sqrt D)
\]
has integer coordinates, norm one, and $1\le\gamma<\varepsilon$.
If $\gamma>1$, both its norm-one coordinates are positive, and its first
coordinate is smaller than $u$, contrary to the definition of $(u,v)$.
Indeed, a first coordinate at least $u$ would force its positive second
coordinate to be at least $v$ by the norm equation, and hence force
$\gamma\geq\varepsilon$.
The decidable alternative is $\gamma=1$, yielding $(x,y)=(x_m,y_m)$.
The strict increase of $x_m$ proves uniqueness.  The index cannot be zero
because $y>0$.
\end{proof}

\subsection{St\o rmer's key divisibility lemma}

\begin{lemma}[Smooth second coordinate forces the fundamental solution]
\label{lem:stormer-key}
Suppose $x^2-Dy^2=1$ with $x,y>0$ and every prime divisor of $y$ divides $D$.
Then $(x,y)$ is the fundamental positive Pell solution for $D$.
\end{lemma}

\begin{proof}
By Lemma~\ref{lem:stormer-pell-powers}, write $(x,y)=(x_m,y_m)$.
We show $m=1$.  Suppose $m>1$ and choose a prime $\ell\mid m$, by bounded
trial division.  The binomial expansion of
$(x_\ell+y_\ell\sqrt D)^{m/\ell}$ shows that $y_\ell\mid y_m$.
Another binomial expansion gives $y_\ell=vC_\ell$, where
\begin{equation}\label{eq:stormer-binomial}
 C_\ell=\sum_{j=0}^{\lfloor(\ell-1)/2\rfloor}
  {\ell\choose 2j+1}u^{\ell-2j-1}v^{2j}D^j.
\end{equation}
Thus every prime dividing $C_\ell$ divides $D$.
The Pell identity gives $\gcd(u,D)=1$.
For any prime $r\mid C_\ell$, reduction of
\eqref{eq:stormer-binomial} modulo $r\mid D$ gives
\[
                   0\equiv\ell u^{\ell-1}\pmod r.
\]
Therefore $r=\ell$.  Since $C_\ell>1$, unique factorization yields
$C_\ell=\ell^e$ for some $e\ge1$ and also $\ell\mid D$.
We now examine the three decidable cases.

If $\ell=2$, then $C_2=2u$.  As $u>1$, a prime divisor of $u$ divides $C_2$
and hence $D$, contradicting $\gcd(u,D)=1$.

If $\ell=3$, the Pell identity reduces the expression to
\[
 C_3=3u^2+Dv^2=4u^2-1=(2u-1)(2u+1).
\]
Both factors exceed one, both are odd, and their greatest common divisor is one
because their difference is two.  They therefore cannot both be powers of
three, contradicting $C_3=3^e$.

Finally let $\ell\ge5$.  In \eqref{eq:stormer-binomial}, every term after the
first is divisible by $\ell^2$.  Indeed, the $j=1$ term contains both $D$ and
${\ell\choose3}$, each divisible by $\ell$; every term with $j\ge2$ contains
$D^2$.  Thus
\[
              C_\ell\equiv\ell u^{\ell-1}\pmod{\ell^2}.
\]
Since $\ell\mid D$ and $\gcd(u,D)=1$, the right-hand side is not divisible
by $\ell^2$.  Hence $e=1$ and $C_\ell=\ell$.
But the first, positive summand of \eqref{eq:stormer-binomial} is
$\ell u^{\ell-1}>\ell$, a contradiction.  Each possible prime divisor of
$m>1$ has been excluded, so $m=1$.
\end{proof}

Notice the special treatment of $\ell=3$.  A claim that the same valuation-one
argument works for every odd prime would be false: the binomial coefficient
${3\choose3}$ contributes no extra factor of three.

\subsection{Reduction of smooth neighbors and the finite algorithm}

\begin{proof}[Proof of Theorem~\ref{thm:stormer-pr}]
If $2\notin S$, every $S$-smooth number is odd, so there are no consecutive
pairs.  Assume henceforth that $2\in S$ and that $n,n+1$ are $S$-smooth.
Factor
\[
             (2n+1)^2-1=4n(n+1)=\prod_{i=1}^k p_i^{a_i}.
\]
Define
\[
 e_i=\begin{cases}
 0,&a_i=0,\\
 1,&a_i\text{ is odd},\\
 2,&a_i>0\text{ is even},
 \end{cases}
 \qquad
 D=\prod_i p_i^{e_i},\qquad
 y=\prod_i p_i^{(a_i-e_i)/2}.
\]
Each exponent is a nonnegative integer.  We have
\[
 (2n+1)^2-Dy^2=1,
\]
and every prime dividing $y$ divides $D$, since positive $a_i$ always gives
positive $e_i$.  The coefficient $D$ is nonsquare: if $D=d^2$, then
\[
 (2n+1-dy)(2n+1+dy)=1
\]
would force $dy=0$, impossible.  Lemma~\ref{lem:stormer-key} now says that
$(2n+1,y)$ is the fundamental solution for $D$.

There are exactly $3^k$ possible coefficient vectors $(e_1,\ldots,e_k)$.
Exactly $2^k$ give squares, namely those with every $e_i\in\{0,2\}$.
A nonsquare coefficient has only one fundamental solution and therefore yields
at most one neighbor pair.  This proves the bound $3^k-2^k$.
Also $D\le R^2$ and $B$ is increasing, so
\[
                 2n+1<B(D)\le B(R^2).
\]

Here is a complete algorithm.  Enumerate all vectors in $\{0,1,2\}^k$ and
form their coefficients $D$.  Discard squares.  For each remaining $D$, search
$1\le y<B(D)$ until $x=\operatorname{isqrt}(Dy^2+1)$ satisfies
$x^2-Dy^2=1$.  The preceding bounded Pell lemma guarantees success.
If $x$ is odd, set $n=(x-1)/2$ and retain it precisely when both $n$ and $n+1$
are positive and $S$-smooth.  Remove duplicate entries and sort.
The construction just proved shows that every required pair occurs;
the final tests show that every returned pair is correct.  The loops are bounded
by explicitly primitive recursive functions.  No stopping test for an infinite
search is being assumed.
\end{proof}

For $S=\{2,3,5\}$ this returns
\[
 (1,2),\ (2,3),\ (3,4),\ (4,5),\ (5,6),\ (8,9),\ (9,10),
 \ (15,16),\ (24,25),\ (80,81).
\]
The accompanying program can compute such small examples while retaining the
proved bound as a totality bound.  The large value of $B(D)$ affects the worst
case bound, not the early return when the fundamental solution is small.

\subsection{Why the proof is available in iPRA}

The mathematical content above reduces to a uniform assertion about finite
codes.  Trial division, primality testing, integer square root, signed integer
arithmetic, exponentiation, binomial coefficients, finite list coding, and
bounded minimization are primitive recursive.  A $D$-smoothness assertion is a
bounded divisibility test.  Algebraic comparisons in $\mathbb Q(\sqrt D)$ are
decidable by integer squaring and sign cases.  The two pigeonhole applications
are finite counting identities.  The power-sequence identities follow by
quantifier-free induction on the index; binomial divisibility follows by the
usual integer product identities and Euclid's lemma.  The apparent universal
properties of lists are bounded predicates and can be replaced by their
primitive recursive Boolean characteristic functions.

For complete proof-theoretic precision, the iPRA assertion can be stated as a
schema for fixed choices of these primitive recursive symbols, with free
natural-number inputs coding $S$ and $n$:
\[
 \operatorname{PrimeList}(S)\ \Longrightarrow\
 \Bigl[\bigl(\operatorname{Smooth}(S,n)\wedge
       \operatorname{Smooth}(S,n+1)\bigr)
 \ \Longleftrightarrow\ n\in\operatorname{Neighbors}(S)\Bigr].
\]
The membership predicate on the right is bounded in the explicitly constructed
output list.  The auxiliary cardinality and height statements are likewise
finite arithmetic statements.  A fully verified implementation must still
supply the iPRA derivations of the lemmas; an ordinary Lean proof over natural
numbers alone is not automatically such a derivation.

\subsection{Lean implementation path}

An implementation can be divided into the following independently checkable
parts.  First encode ordered quadratic pairs and prove the integer sign test,
ring identities, norm identities, and the formulas used by the quotient
construction.  Next formalize the two bounded pigeonhole searches and prove the
output specification of the Pell constructor.  Then prove generation by the
fundamental solution, the displayed binomial formula, and the three prime-index
cases.  Finally define the coefficient enumeration and smoothness filters,
and prove soundness, completeness, the cardinal bound, and the explicit height
bound.  Every unbounded-looking existential in this route has already been
replaced above by a stated finite search bound.

For an iPRA certificate rather than only a metatheoretic Lean proof, use the
repository's primitive recursive syntax and derivation infrastructure to
represent those definitions, and translate the displayed bounded induction
invariants into characteristic-function equations.  The substantive remaining
formalization work is the checked translation and arithmetic infrastructure;
there is no omitted Diophantine estimate or primitive-divisor theorem in this
St\o rmer argument.
